\begin{ThreePartTable}

\begin{TableNotes}

    \item[1] "SHA-256" 表示 SHA-256 算法，参考 \ref{mod:sha} \textit{\nameref{mod:sha}} 章节。

\end{TableNotes}

\begin{longtable}[c]{ |c|c|c|c|c| }
\caption{根据 \begin{math}Key_{A}\end{math}\begin{math}、Key_{B}\end{math} 生成的 \begin{math}Key\end{math} 值}
\label{tab:extmemencr-key-possibility}
\\\hline 
\rowcolor{lightgray}
\hyperref[fielddesc:EFUSEXTSKEYLENGTH256]{EFUSE\_XTS\_KEY\_LENGTH\_256} & \begin{math}Key\end{math}  & \begin{math}Key\end{math} 长度（位） \\\hline
\endfirsthead
\hline
\rowcolor{lightgray}
\hyperref[fielddesc:EFUSEXTSKEYLENGTH256]{EFUSE\_XTS\_KEY\_LENGTH\_256} & \begin{math}Key\end{math}  & \begin{math}Key\end{math} 长度（位） \\\hline
\endhead
\insertTableNotes
\endfoot
1       &   \begin{math}           \{Key_{B},Key_{A}\}\end{math}    &   256   \\\hline
0        &   \begin{math}           SHA-256(Key_{A})\end{math}\tnote{1}    &   256   \\\hline
\end{longtable}
\end{ThreePartTable}
